\begin{abstract}
In the domain of membership query problems---determining whether a given element belongs to a predefined set---Hash Tables and Bloom Filters represent a classical space--accuracy trade-off. Hash Tables provide exact, deterministic lookups at the cost of significant per-element memory overhead, whereas Bloom Filters offer a compact probabilistic alternative that permits a controllable false positive rate (FPR). This paper presents a systematic empirical study that identifies the \emph{Performance Crossover Point}: the dataset scale $N$ at which a Hash Table, constrained by a hard memory budget $M$, begins to degrade catastrophically while a Bloom Filter continues to operate with stable throughput. We implement both data structures from scratch in Java---a Hash Table with Separate Chaining and a Bloom Filter with Kirsch--Mitzenmacher double hashing ($k=7$, target FPR $p=1\%$)---and benchmark them across dataset sizes from $N = 1{,}000$ to $N = 2{,}000{,}000$ under five JVM memory caps ($16$--$256$~MB). Our results demonstrate that the Hash Table consumes approximately $100$~bytes per element versus $1.2$~bytes for the Bloom Filter (an ${\approx}83\times$ compression ratio), and that JVM garbage collection pressure causes Hash Table throughput to collapse shortly before an \texttt{OutOfMemoryError}. The empirical FPR of the Bloom Filter remains tightly bounded within $0.84\%$--$1.38\%$, closely matching the Mitzenmacher theoretical prediction. We provide concrete crossover thresholds for each memory tier and discuss practical implications for system design in CDN caching, database pre-filtering, and network deduplication.
\end{abstract}
